\documentclass[]{article}

\usepackage{amsmath}
\usepackage{multirow}
\usepackage{natbib}
\usepackage{graphicx} 
\usepackage[margin=1in]{geometry}

\begin{document}

In this study, we investigated the nature of anomalous diffusion and the evolution of Lévy flight statistics in two-dimensional turbulence undergoing an inverse energy cascade. The primary challenge addressed was the non-stationary nature of Lagrangian transport, which complicates the application of standard continuous-time random walk models that typically assume stationary statistics.

To explore this, we utilized velocity fields generated from 1024$^2$ direct numerical simulations and employed a wavelet-based filtering approach to isolate eddy dynamics across multiple scales. We quantified the Lévy stability exponent $\alpha$ using the McCulloch quantile method and evaluated its temporal evolution alongside the spectral condensation process. Furthermore, we analyzed the influence of Lagrangian coherent structures, such as hyperbolic points, on the observed transport statistics.

Our results demonstrate that the stability exponent $\alpha$ is both non-stationary and scale-dependent, decreasing from approximately 2.0 at smaller scales to 1.75 at the largest scales as the energy spectrum condenses. Statistical tests confirmed the non-stationarity of $\alpha(t)$, and our comparison with continuous-time random walk predictions revealed a significant discrepancy, indicating that the relationship between eddy lifetimes and Lévy statistics is more complex than simple theoretical frameworks suggest. Notably, our causal analysis indicates that the observed drift in $\alpha$ is not directly driven by the instantaneous spectral peak or the spatial density of hyperbolic points.

We conclude that the complexity of Lagrangian dynamics in two-dimensional turbulence arises from intricate, multi-scale interactions that are not fully captured by instantaneous spectral evolution or local structural classifications. These findings suggest that a more comprehensive understanding of transport in turbulent flows requires moving beyond simple random walk models to account for the non-local and multi-scale nature of the underlying velocity fields.

\end{document}
                